\begin{center}
  \begin{tikzpicture}[]
    \newcommand{\spacing}[0]{6mm}
    \newcommand{\projectWidth}[0]{50mm}
    \newcommand{\projectHeight}[0]{28mm}
    \newcommand{\bgcolor}{white}
    
    \tikzstyle{arrow} = [thick,->,>=stealth, draw=black]
    \tikzstyle{box} = [rectangle,thick,draw=black]
    \tikzstyle{project} = [
      box,
      fill=black!5!\bgcolor,
      draw=\bgcolor!5!black,
      minimum height=\projectHeight,
      minimum width=\projectWidth,
      text width=\projectWidth,
    ]
    \tikzstyle{file} = [
      box,
      fill=black!10!\bgcolor,
      draw=\bgcolor!10!black,
      minimum height=(\projectHeight-3*\spacing)/2,
      minimum width=\projectWidth-2*\spacing,
      text width=\projectWidth-2*\spacing
    ]
    
    \coordinate (origin) at (0,0);
    
    \node[project,anchor=east] (codeproject) at ([xshift=-\spacing] origin) {};
    \node[anchor=south west] () at (codeproject.north west) {\small \filename{Codebase}};
    \node[file,anchor=south] (codecs)  at ([yshift= \spacing/2] codeproject) {\filename{Program.cs}};
    \node[file,anchor=north] (codeprj) at ([yshift=-\spacing/2] codeproject) {\filename{Codebase.csproj}};
    
    \node[project,anchor=west] (testproject) at ([xshift= \spacing] origin) {};
    \node[anchor=south west] () at (testproject.north west) {\small \filename{UnitTests}};
    \node[file,anchor=south] (testcs)  at ([yshift= \spacing/2] testproject) {\filename{UnitTest1.cs}};
    \node[file,anchor=north] (testprj) at ([yshift=-\spacing/2] testproject) {\filename{UnitTests.csproj}};
    
    \draw[arrow] (testprj) -- (codeprj);
  \end{tikzpicture}
\end{center}
